\begin{tikzpicture}[
  level/.style={draw, thick, minimum width=3.5cm, align=center},
  >={Stealth[scale=1.2]},
  pump/.style={->, thick, red!70!black},
  stim/.style={->, thick, blue!70!black, dashed},
  spon/.style={->, thick, decorate, decoration={snake, amplitude=2pt, segment length=6pt}, blue!50!black}
]

% --- Położenie poziomów energetycznych (oś Y) ---
\coordinate (G) at (0,0);   % Stan podstawowy
\coordinate (M) at (0,3.5); % Stan metastabilny
\coordinate (E) at (0,7);   % Stan wzbudzony

% --- Rysowanie poziomów ---
\draw[level] (-2.2,0) -- (2.2,0) node[right, xshift=0.8cm] {\(1^3S\)};
\draw[level] (-2.2,3.5) -- (2.2,3.5) node[right, xshift=0.8cm] {\(2^3S\) (metastable) \\ \(\tau = 1142\) ns};
\draw[level] (-2.2,7) -- (2.2,7) node[right, xshift=0.8cm] {\(3^3P\)};

% --- 1. Wzbudzenie laserowe (główna strzałka pompy) ---
\draw[pump] (-3,0.2) -- (-3,6.8) 
    node[midway, left, yshift=-0.6cm, fill=white, inner sep=2pt] {Laser 205 nm};

% --- 2. Rozpad spontaniczny z powrotem do stanu podstawowego (dominujący) ---
\draw[spon] (-0.8,7) -- (-0.8,0) 
    node[midway, left, yshift=0.6cm, fill=white, inner sep=2pt] {Spontaneous (90.3\%)};

% --- 3. Rozpad spontaniczny do stanu 2^3S (pożądany kanał) ---
\draw[spon] (0.8,7) -- (0.8,3.5) 
    node[midway, right, yshift=-0.6cm, fill=white, inner sep=2pt] {Spontaneous (12\%)};

% --- 4. Rozpad wymuszony (stymulowany) do 2^3S - metoda zwiększająca wydajność ---
\draw[stim] (2.6,7) -- (2.6,3.5) 
    node[midway, right, yshift=0.6cm, fill=white, inner sep=2pt] {Stimulated (\(\pi\)-impuls)};

% --- Opis procesu pod schematem ---
% \node at (0, -1.8) { \textbf{Schemat kaskady:} \(1^3S \xrightarrow{205\text{ nm}} 3^3P \xrightarrow{12\%} 2^3S\) };

% --- Dodatknotka o wydajności ---
% \node at (0, -2.8) [text width=10cm, align=center, font=\small] {
    % Uwaga: Stosując technikę STIRAP lub szybkie wygaszanie pola laserowego, 
    % można zwiększyć populację stanu \(2^3S\) kosztem rozpadu do \(1^3S\).
% };

\end{tikzpicture}
% \end{document}